% This is file `memo-symbols-standard.tex'.
% This file is part of GraphTheorySymbol.
%
% GraphTheorySymbol is free software:
% you can redistribute it and/or modify it
% under the terms of the GNU Lesser General Public License
% as published by the Free Software Foundation,
% either version 3 of the License,
% or (at your option) any later version.
%
% GraphTheorySymbol is distributed
% in the hope that it will be useful,
% but WITHOUT ANY WARRANTY;
% without even the implied warranty of
% MERCHANTABILITY or FITNESS FOR A PARTICULAR PURPOSE.
% See the GNU Lesser General Public License for more details.
%
% You should have received a copy of
% the GNU Lesser General Public License
% along with GraphTheorySymbol.
% If not, see <https://www.gnu.org/licenses/>.
%
% This work is maintained by Laurent Lyaudet.
%
% This work consists of the files graph-theory-symbol.dtx,
%                                 graph-theory-symbol.ins,
%                                 tex/memo-symbols-standard.tex,
%                                 tex/memo-symbols-monochrome.tex,
%                                 tex/memo-symbols-meta-logic.tex,
%                                 tex/memo-preamble.tex,
%                                 tex/memo-versatile-macros.tex,
%                                 graph-theory-symbol-memo.tex, and
%                                 test.tex,
% and the derived files           graph-theory-symbol.sty,
%                                 graph-theory-symbol-doc-A4.pdf,
%                                 graph-theory-symbol-doc-USLetter.pdf,
%                                 graph-theory-symbol-doc-Book.pdf,
%                                 graph-theory-symbol-memo-A4.pdf,
%                                 graph-theory-symbol-memo-USLetter.pdf,
%                                 graph-theory-symbol-memo-Book.pdf,
%                                 test-A4.pdf,
%                                 test-USLetter.pdf, and
%                                 test-Book.pdf.
%
% ©Copyright 2026 Laurent Frédéric Bernard François Lyaudet

\begin{tabular}{@{}llll@{}}
Long name & Short name & Result & In context\\
|\GTSmetaNot/| & |\GTSmN/| & \GTSmN/ & \GTSmN/\GTSa/\\
|\GTSmetaAnd/| & |\GTSmA/| & \GTSmA/
  & \GTSxa[mC=red]\GTSmA/\GTSxa[mC=green]\\
|\GTSmetaOr/| & |\GTSmO/| & \GTSmO/
  & \GTSxa[mC=red]\GTSmO/\GTSxa[mC=green]\\
\end{tabular}

Meta-logic can be used for compound adjacency types, etc.

Implicit meta-and via concatenation,
no space except around whole relations compositions,
no parentheses via higher priority to meta-and, brackets when all possible
adjacencies are specified between two vertices:\\
Tournament adjacency:\\
$\forall u,v,
u \mathrel{[\GTSdAD/\GTSnDAU/\GTSmO/\GTSnDAD/\GTSdAU/]} v$.\\
Tournament adjacency with maybe some colored edge also, who knows?:\\
$\forall u,v,
u \mathrel{\GTSdAD/\GTSnDAU/\GTSmO/\GTSnDAD/\GTSdAU/} v$.\\
Directed graph:\\
$\forall u,v,
u \mathrel{[\GTSdAD/\GTSnDAU/\GTSmO/\GTSnDAD/\GTSdAU/
\GTSmO/\GTSdAD/\GTSdAU/\GTSmO/\GTSnDAD/\GTSnDAU/]} v$.\\
Oriented graph:\\
$\forall u,v,
u \mathrel{[\GTSdAD/\GTSnDAU/\GTSmO/\GTSnDAD/\GTSdAU/
\GTSmO/\GTSnDAD/\GTSnDAU/]} v$.\\
Tournament adjacency:\\
$\mathrel{[\GTSdAD/\GTSnDAU/\GTSmO/\GTSnDAD/\GTSdAU/]}$
is obtained with |\mathrel{[\GTSdAD/\GTSnDAU/\GTSmO/\GTSnDAD/\GTSdAU/]}|.\\
Tournament adjacency with maybe some colored edge also, who knows?:\\
$\mathrel{\GTSdAD/\GTSnDAU/\GTSmO/\GTSnDAD/\GTSdAU/}$
is obtained with |\mathrel{\GTSdAD/\GTSnDAU/\GTSmO/\GTSnDAD/\GTSdAU/}|.\\
Directed graph:\\
$\mathrel{[\GTSdAD/\GTSnDAU/\GTSmO/\GTSnDAD/\GTSdAU/
\GTSmO/\GTSdAD/\GTSdAU/\GTSmO/\GTSnDAD/\GTSnDAU/]}$ is obtained with\\
|\mathrel{[\GTSdAD/\GTSnDAU/\GTSmO/\GTSnDAD/\GTSdAU/|\\
|\GTSmO/\GTSdAD/\GTSdAU/\GTSmO/\GTSnDAD/\GTSnDAU/]}|.\\
Oriented graph:\\
$\mathrel{[\GTSdAD/\GTSnDAU/\GTSmO/\GTSnDAD/\GTSdAU/
\GTSmO/\GTSnDAD/\GTSnDAU/]}$ is obtained with\\
|\mathrel{[\GTSdAD/\GTSnDAU/\GTSmO/\GTSnDAD/\GTSdAU/|%
|\GTSmO/\GTSnDAD/\GTSnDAU/]}|.

Prefix notation:\\
$\exists u,v, u
\mathrel{[\GTSmO/\GTSxa[mC=red,d]\GTSxa[mC=green,d]\GTSxa[mC=blue,u]]}
v$.\\
$\exists v,e, v
\mathrel{[\GTSmO/\GTSxi[mC=red,i]\GTSxi[mC=green,i]\GTSxi[mC=blue,o]]}
e$.\\
$\mathrel{[\GTSmO/\GTSxa[mC=red,d]\GTSxa[mC=green,d]\GTSxa[mC=blue,u]]}$
is obtained with\\
|\mathrel{[\GTSmO/\GTSxa[mC=red,d]\GTSxa[mC=green,d]\GTSxa[mC=blue,u]]}|.\\
$\mathrel{[\GTSmO/\GTSxi[mC=red,i]\GTSxi[mC=green,i]\GTSxi[mC=blue,o]]}$
is obtained with\\
|\mathrel{[\GTSmO/\GTSxi[mC=red,i]\GTSxi[mC=green,i]\GTSxi[mC=blue,o]]}|.
